\section{Метрики и протокол агрегирования}\label{sec:metrics}

Основной целевой метрикой остаётся социальное благосостояние
\begin{equation}
SW = \sum_{t}\sum_{i}\sum_{j} x_{i,j}(t)\,\bigl(U_i(t,j)-C_j(t,i)\bigr),
\end{equation}
то есть суммарный чистый вклад принятых задач в систему с учётом полезности
и стоимости исполнения. Именно эта метрика непосредственно согласована с
математической постановкой главы~2 и с логикой аукционного слоя MA--VCG.

Наряду с $SW$ в экспериментах использовались следующие эксплуатационные
метрики:
\begin{itemize}
	\item \textbf{acceptance rate} --- доля задач, принятых к исполнению;
	\item \textbf{deadline success rate} --- доля завершённых задач, уложившихся
	в дедлайн;
	\item \textbf{average latency} --- средняя сквозная задержка выполнения;
	\item \textbf{fairness index} --- индекс Джайна по реализованному
	распределению нагрузки;
	\item \textbf{gini payment} --- коэффициент Джини по VCG-платежам;
	\item \textbf{load imbalance} --- дисбаланс загрузки узлов;
	\item \textbf{completed tasks} и \textbf{completed before deadline} ---
	характеристики фактической пропускной способности.
\end{itemize}

Для обучаемых методов дополнительно фиксировались финальные значения
\texttt{training social welfare}, \texttt{training reward}, TD-ошибки и
$\varepsilon$-параметра, а также полные кривые обучения по эпизодам.

Агрегирование результатов выполнялось по пяти независимым seed. В таблицах
главы приводятся значения <<mean $\pm$ std>>, а в агрегированных CSV-файлах
дополнительно сохраняются доверительные интервалы $95\%$ по всем основным
метрикам. Такая схема не заменяет полноценные парные статистические тесты,
но уже позволяет отделять устойчивые различия от случайных колебаний между
прогонами.

Отдельно следует отметить специфику платежной fairness-метрики. На текущем
валидационном стенде коэффициент Джини по VCG-платежам оказался близок к
нулю почти во всех сценариях. Это означает, что при сравнительно разреженном
потоке задач и умеренном числе конкурирующих заявок внешний эффект
отдельного устройства на остальных участников часто мал, а потому сам
платёжный сигнал слабо различает методы. По этой причине интерпретация
справедливости в настоящей главе опирается в первую очередь на
\texttt{fairness index}, пошаговую динамику и тепловые карты загрузки узлов.
